%%=============================================================================
%% Voorwoord
%%=============================================================================

\chapter*{\IfLanguageName{dutch}{Woord vooraf}{Preface}}%
\label{ch:voorwoord}

%% TODO:
%% Het voorwoord is het enige deel van de bachelorproef waar je vanuit je
%% eigen standpunt (``ik-vorm'') mag schrijven. Je kan hier bv. motiveren
%% waarom jij het onderwerp wil bespreken.
%% Vergeet ook niet te bedanken wie je geholpen/gesteund/... heeft

Mijn interesse voor dit onderwerp is ontstaan vanuit een fascinatie voor zowel radiotechnologie
als netwerktechnologie, en in het bijzonder voor de achterliggende technieken voor dataoverdracht
en de benodigde infrastructuur. Bij het bestuderen van communicatie-infrastructuur groeide het
besef dat deze infrastructuur ook zijn grenzen kent en in bepaalde situaties verrassend fragiel
kan zijn --- denk aan afgelegen gebieden, maritieme omgevingen of rampscenario's waarbij de
bestaande netwerken wegvallen. Die kwetsbaarheid wekte de vraag op of radioverbindingen een
waardig alternatief kunnen vormen voor conventionele netwerkinfrastructuur, en vormde zo de
aanleiding voor deze bachelorproef.

Graag wil ik een aantal mensen bedanken die mij tijdens dit traject hebben ondersteund.
In de eerste plaats mijn co-promotor Frans Van de Velde, wiens praktijkervaring en technische
inzichten een waardevolle leidraad vormden doorheen het onderzoek. Verder dank ik mijn
promotor Peter Beke voor de begeleiding en de constructieve feedback tijdens het uitwerken
van dit voorstel. Tot slot een bijzonder woord van dank aan mijn lokale UBA-sectie, die steeds
bereid was om goede raad te geven over radiotechnische vraagstukken en nieuwe ideeën aan
te reiken die het onderzoek verder hebben verrijkt.